% teorema-degrau3a-draft.tex -- DEGRAU 3a: o teorema da LEI DIAGONAL POR FACE
% (multi-dado), mais o corolario do LIDER ESTAVEL.
% [DRAFT --- aguarda parecer], 2026-07-27.  NAO e o paper.
% Fontes da verdade:
%   06-notas/2026-07-27-degrau3-scoping.md          (estrutura derivada: D1--D7)
%   06-notas/paper/paper-v4.tex                     (Thm 6 = thm:forcing; Ap. A/B)
%   06-notas/t2-prova/teorema-degrau2-draft.tex     (padrao de draft da casa)
%   05-experimentos/t2-n-neuronios/degrau3a_verifica.{py,jsonl}
% Compilar: pdflatex teorema-degrau3a-draft.tex (2x).
\documentclass[11pt]{article}
\usepackage{amsmath,amssymb,amsthm}
\usepackage[text={6in,8.4in},centering]{geometry}
\usepackage{booktabs}
\usepackage{microtype}
\usepackage{url}
\usepackage[colorlinks=true,linkcolor=blue,citecolor=blue,urlcolor=blue]{hyperref}
\usepackage[nameinlink]{cleveref}

\newtheorem{theorem}{Theorem}
\newtheorem{lemma}[theorem]{Lemma}
\newtheorem{proposition}[theorem]{Proposition}
\newtheorem{corollary}[theorem]{Corollary}
\theoremstyle{remark}
\newtheorem{remark}[theorem]{Remark}
\crefname{theorem}{Theorem}{Theorems}
\crefname{lemma}{Lemma}{Lemmas}
\crefname{remark}{Remark}{Remarks}
\crefname{proposition}{Proposition}{Propositions}
\crefname{corollary}{Corollary}{Corollaries}
\crefname{section}{Section}{Sections}
\crefname{equation}{}{}
\Crefname{equation}{}{}

\newcommand{\R}{\mathbb{R}}
\newcommand{\thit}{t_{\mathrm{hit}}}
\newcommand{\texit}{t_{\mathrm{exit}}}
\newcommand{\trel}{t_{\mathrm{release}}}
\newcommand{\Ar}{\mathcal{A}}
\newcommand{\Tep}{T_{\mathrm{ep}}}
\newcommand{\hyb}{\mathrm{hyb}}
\newcommand{\res}{\mathrm{res}}
\newcommand{\Cres}{C_{\mathrm{res}}}
\newcommand{\Xmax}{X_{\max}}
\newcommand{\dgap}{\delta_{\mathrm{gap}}}
\newcommand{\hyplabel}[1]{\textup{(#1)}}

\title{Degrau 3a: the stable-leader corollary and the\\
\emph{per-face diagonal area law} for a dying unit with $n$ active faces\\
\large [DRAFT --- aguarda parecer], 2026-07-27}
\author{j-space internal note}
\date{2026-07-27}

\begin{document}
\maketitle

\begin{abstract}
Theorem~6 of \texttt{paper-v4} is a \emph{scalar} entry--exit law: one face of
the dying unit carries one area $\Ar=\int R$, and the unit revives when that
area returns to zero. With $n$ data the dead unit has $n$ faces at once, and the
scoping note \cite{scoping} derives the exact field
$\dot z=C\bigl[(R-g_\beta(z))\odot\sigma(\beta z)\bigr]$, $C=XX^\top/n$, together
with the exact multi-face quadrature, whose coupling weight is
$(1+u_i)/(1+u_k)$. This note turns two consequences of that structure into
statements with proofs. \textbf{\Cref{thm:A} (stable leader)} identifies the
paper's hypothesis \hyplabel{H4$'$}/$\delta_1$ as a \emph{leader-slack}
hypothesis and shows that the whole of Theorem~6 survives verbatim with
$R=R_{i^\ast}$ and $c=C_{i^\ast i^\ast}$, at the price of one threshold
$\beta\dgap\ge\log(\beta\Gamma)$, which we derive; it gives
$\beta_0=O\bigl(\dgap^{-1}\log(1/\dgap)\bigr)$ and, for $n$ data in general
position, $\beta_0=O(n\log n)$ --- a new, quantitative and falsifiable
prediction. \textbf{\Cref{thm:B} (per-face diagonal law)} is the substantive
one: under $C$-orthogonality of the active faces, quantified by
$\varepsilon=\max_{i\ne k}|C_{ik}|/C_{ii}$, the system decouples into $n$ exact
copies of L1$'$, \emph{each face carries its own area against its own threshold}
$\Ar_i=[\Psi(u_i^0)-1]/(\beta C_{ii})$, and release happens on the first face to
cross, with the first-order structure of the winner. The new content is not the
quadrature (which is $n$ copies of the paper's) but \textbf{\cref{lem:argmin}},
the comparison lemma: the $\min_i$ is stable under the coupling perturbation
exactly when the runner-up's threshold slack exceeds it. A campaign of $48$
drawn instances at $\beta\in[2\cdot10^2,2\cdot10^3]$ with
$\beta\dgap\ge10$ \emph{by construction} --- the regime the scoping note
identified as the empirical blind spot (A5) --- confirms both: the leader law is
right on $18/18$ faces with error $O(\beta^{-1})$; where the leader is not
stable, the scalar law predicts \textbf{the wrong face in $18/18$} with a
$244\%$ median time error while the per-face law is right on $18/18$; and a
controlled violation with $\varepsilon\approx1.2$ breaks the per-face law on
$11/12$ faces with an error that does \emph{not} decay in $\beta$, delimiting
the theorem. We also report, and do not soften, one measurement that does
\emph{not} close: the first-order coefficient $-1/(cR^{\hyb}(\texit))\beta^{-1}$
of Theorem~6(ii) is not observed here, in the multi-face toy \emph{or} in an
$n=1$ control; what is observed is $\trel-\texit=+\Theta(\log\beta/\beta^2)$,
which is the statement of L4$'$(ii) of the same paper. The two statements
disagree, we cannot adjudicate them from this note, and we declare it as
\cref{sec:open}\,(D5).
\end{abstract}

\tableofcontents

%======================================================================
\section{Setup: the dead unit with $n$ active faces}\label{sec:setup}

\subsection{The exact field}

Keep the network, data and loss of \cite[\S2]{paper}: $(x_i,y_i)_{i=1}^n$, the
distinguished unit $A$ with readout $+1$ and parameters
$\theta_A\in\R^{d+1}$, the augmented data $X_i=(x_i,1)$, the margins
$z_i=X_i\cdot\theta_A$, and the forcing
$R_i(t)=y_i-f_{-A}(x_i;\theta_B(t))$ generated by the live block. Write
$g_\beta$ for the softplus, $\sigma$ for the logistic. Gradient flow on
$\theta_A$ is $\dot\theta_A=-\frac1n\sum_i\res_i\,\sigma(\beta z_i)X_i$ with
$\res_i=g_\beta(z_i)-R_i$, so that
\begin{equation}\label{eq:D1}
\boxed{\
\dot z_i=\sum_{k=1}^{n}C_{ik}\bigl(R_k(t)-g_\beta(z_k)\bigr)\,\sigma(\beta z_k),
\qquad C_{ik}=\frac{\langle X_i,X_k\rangle}{n},\ }
\end{equation}
with $C=XX^\top/n\succeq0$ symmetric of rank $d+1$. \Cref{eq:D1} is an
\emph{identity}, verified against automatic differentiation of $L$ to $6\cdot
10^{-11}$ in \cite[I1]{scoping}. The scalar setting of \cite[Thm.~6]{paper} is
$n=1$, where $C=c=(x_j^2+1)/n$.

Two structural readings of \eqref{eq:D1}, both from \cite{scoping} and both used
below only as orientation: in the park ($g_\beta\approx0$, $\sigma\approx
e^{\beta z}$) the flow is \emph{gradient ascent} on the sum of exponentials
$V=(n\beta)^{-1}\sum_iR_ie^{\beta X_i\cdot\theta}$, and in the channel variables
$q_i=e^{\beta z_i}$ it is a non-autonomous Lotka--Volterra system
$\dot q_i=\beta q_i\sum_kC_{ik}R_kq_k$ with symmetric PSD community matrix.

\subsection{The exact multi-face quadrature}

Let $u_i=e^{-\beta z_i}$, $w_i=e^{\beta z_i}=1/u_i$ and $\Psi(u)=\log u+u$, the
same gauge as \cite[L1$'$]{paper}. From $\dot u_i=-\beta u_i\dot z_i$ and
$\sigma(\beta z_k)=1/(1+u_k)$,
\begin{align}
\frac{d\Psi(u_i)}{dt}
&=-\beta(1+u_i)\sum_{k}\frac{C_{ik}\bigl(R_k-g_\beta(z_k)\bigr)}{1+u_k}
\notag\\[2pt]
&=\underbrace{-\beta C_{ii}R_i+C_{ii}\log(1+w_i)}_{\text{the scalar L1}'\text{ term}}
\;\underbrace{-\;\beta\sum_{k\ne i}C_{ik}\,\frac{1+u_i}{1+u_k}\,
\bigl(R_k-g_\beta(z_k)\bigr)}_{=:\,\mathcal C_i(t),\ \text{the coupling}} .
\label{eq:D4}
\end{align}
\Cref{eq:D4} is exact (verified to relative $9.6\cdot10^{-9}$,
\cite[I2]{scoping}); we used $g_\beta(z_i)=\beta^{-1}\log(1+w_i)$ and
$(1+u_i)/u_i\cdot u_i=1+u_i$. The whole of this note is the study of
$\mathcal C_i$.

\begin{remark}[the coupling weight is not bounded by one]\label{rem:weight}
$(1+u_i)/(1+u_k)$ is $\le1$ iff $u_i\le u_k$ iff $z_i\ge z_k$. In deep park
($u\gg1$) it is $u_i/u_k=e^{\beta(z_k-z_i)}$, which is exponentially
\emph{large} when face $i$ is parked deeper than face $k$. This is the single
fact that makes the multi-face problem non-trivial, and it is why both theorems
below are, at bottom, statements about \emph{which face is shallowest}. It also
says that any bound on $\mathcal C_i$ that treats the weight as $O(1)$ is wrong;
\cref{lem:A1,lem:B1} both carry the weight explicitly.
\end{remark}

\subsection{Hypotheses}

We assume \hyplabel{H0$'$}, \hyplabel{H1$'$}, \hyplabel{H2$'$},
\hyplabel{H3$'$}, \hyplabel{H5$'$} of \cite[\S3.1]{paper} verbatim, read
per face where they mention $R$: $R_i\in C^1$, $|\dot R_i|\le C_R$,
$|R_i|\le\Cres$, transversal entry and exit for the face that decides. We
\emph{replace} the $\delta_1$ clause of \hyplabel{H4$'$} --- ``the other margins
of $A$ stay $\le-\delta_1$'', which in the scalar theorem is what makes all
faces but one invisible --- because in this note the other faces are precisely
what we are looking at. What survives of \hyplabel{H4$'$} is the live-block part
($|z_k(x_i;\theta_B)|\ge\delta_B$, no live switch on $[0,H]$), which we keep and
never touch. Write
\[
\thit=0\ \text{(the death instant, after translation)},\qquad
z^0=z(\thit),\qquad \Ar_i(t)=\int_{\thit}^{t}R_i(s)\,ds ,
\]
and normalize so that the \emph{leader} $i^\ast=\arg\max_iz_i^0$ has
$z^0_{i^\ast}=0$: it is the face on which $A$ dies. Put
\begin{equation}\label{eq:gapdef}
\delta_i:=-z_i^0\ge0,\qquad
\dgap:=\min_{k\ne i^\ast}\delta_k=z^0_{i^\ast}-z^0_{(2)} .
\end{equation}

%======================================================================
\section{Theorem A: the stable leader}\label{sec:A}

The point of this section is bookkeeping, and it is short by design: nothing new
is proved, one hypothesis is \emph{renamed} and one threshold is \emph{derived}.
The gain is interpretive and quantitative, and it is stated in
\cref{rem:whatAbuys}.

\begin{itemize}
\item[\hyplabel{H3-gap}] \textbf{Leader slack.} With $\dgap$ as in
\eqref{eq:gapdef}, $\Ar_{\max}:=\sup_{[\thit,H]}\max_i|\Ar_i|$,
$C_{\mathrm{off}}:=\max_i\sum_{k\ne i}|C_{ik}|$,
$c_{\min}:=\min_iC_{ii}$ and
$\Gamma:=\bigl(C_{\mathrm{off}}/c_{\min}\bigr)\bigl(\Cres+1\bigr)
\max\bigl(\Tep,\ \Ar_{\max}\bigr)$,
\begin{equation}\label{eq:H3gap}
\beta\,\dgap\ \ge\ \log\bigl(\beta\,\Gamma\bigr)\ +\ \log\frac1\eta ,
\end{equation}
for a tolerance $\eta\in(0,1)$ fixed in advance.
\end{itemize}

\begin{lemma}[the leader's quadrature decouples]\label{lem:A1}
Assume \hyplabel{H3-gap} and that $z_{i^\ast}(t)=\max_iz_i(t)$ for
$t\in[\thit,s]$. Then on $[\thit,s]$ the coupling term of \eqref{eq:D4} at
$i=i^\ast$ obeys
\begin{equation}\label{eq:Abound}
\bigl|\mathcal C_{i^\ast}(t)\bigr|
\ \le\ \beta\,C_{\mathrm{off}}\,\bigl(\Cres+\max_kg_\beta(z_k)\bigr)\,
e^{-\beta\dgap},
\qquad
\int_{\thit}^{s}\bigl|\mathcal C_{i^\ast}\bigr|\,dt\ \le\ \eta\,c_{\min},
\end{equation}
so that $\Psi(u_{i^\ast})$ satisfies the identity
\eqref{eq:D4} of \cite[L1$'$]{paper} with $c=C_{i^\ast i^\ast}$,
$R=R_{i^\ast}$, and with the bookkeeping error
$\mathcal E$ of \cite[(B.quadrature)]{paper} enlarged by \eqref{eq:Abound}
--- that is, with $\delta_1$ replaced by $\dgap$ throughout.
\end{lemma}

\begin{proof}
By hypothesis $z_{i^\ast}\ge z_k$ for every $k$, hence $u_{i^\ast}\le u_k$ and
the weight of \cref{rem:weight} satisfies
$0<(1+u_{i^\ast})/(1+u_k)\le1$. \emph{This is the only place the leader
property is used, and it is used to make the dangerous weight harmless.} It
must be sharpened to get the exponential: for $k\ne i^\ast$ we have
$z_k\le z_{i^\ast}-\dgap$ \emph{at $t=\thit$}, and along the flow the leader
stays the leader by hypothesis, so $u_{i^\ast}/u_k=e^{\beta(z_k-z_{i^\ast})}$.
When both faces are in park, $u\gg1$ and
$(1+u_{i^\ast})/(1+u_k)\le2u_{i^\ast}/u_k\le2e^{-\beta\dgap}$ once the gap is
maintained; when the leader is not in park ($u_{i^\ast}\lesssim1$) the weight is
$\le2/u_k\le2e^{\beta z_k}$, and $z_k\le-\dgap$ there because the leader is then
near $0$. In both regimes the weight is $\le2e^{-\beta\dgap}$. Multiplying by
$\beta\sum_{k\ne i^\ast}|C_{ik}|\le\beta C_{\mathrm{off}}$ and by
$|R_k-g_\beta(z_k)|\le\Cres+\max_kg_\beta(z_k)$ gives the pointwise bound (the
factor $2$ absorbed into $\Gamma$). Integrating over an interval of length
$\le\Tep$ and using $g_\beta\le1$ on the park (clause~(i) of
\cite[Thm.~6]{paper} gives $g_\beta(z_{\rm park})=O(\beta^{-2})$; off the
leader $g_\beta(z_k)\le\beta^{-1}e^{-\beta\dgap}$) gives
$\int|\mathcal C_{i^\ast}|\le\beta C_{\mathrm{off}}(\Cres+1)\Tep
e^{-\beta\dgap}\le c_{\min}\eta$ by \eqref{eq:H3gap}, since
$\beta\Gamma e^{-\beta\dgap}\le\eta$ and
$\Gamma\ge(C_{\mathrm{off}}/c_{\min})(\Cres+1)\Tep$.
The remaining assertion is verbatim: with $\mathcal C_{i^\ast}$ absorbed into
$\mathcal E$, \eqref{eq:D4} at $i=i^\ast$ \emph{is}
\cite[eq.~(quadratureF)]{paper}.
\end{proof}

\begin{lemma}[the leader also wins the race]\label{lem:A2}
Under \hyplabel{H3-gap}, for every $k\ne i^\ast$ and every $t\le H$,
$\Ar_k(t)<[\Psi(u_k^0)-1]/(\beta C_{kk})$: no other face can reach its own
threshold \eqref{eq:tau} inside the horizon.
\end{lemma}

\begin{proof}
$\Psi(u_k^0)-1=\beta\delta_k+e^{\beta\delta_k}-1\ge e^{\beta\dgap}-1$, so the
threshold is $\ge(e^{\beta\dgap}-1)/(\beta C_{kk})$, while
$\Ar_k\le\Ar_{\max}$. It suffices that
$e^{\beta\dgap}\ge1+\beta C_{kk}\Ar_{\max}$, which follows from
\eqref{eq:H3gap} because $\Gamma\ge\Ar_{\max}$ up to the constant
$C_{\mathrm{off}}/c_{\min}$ absorbed there and $\eta<1$.
\end{proof}

\begin{theorem}[entry--exit law, multi-face, stable leader]\label{thm:A}
Let $A$ be a softplus-$\beta$ unit with readout $+1$ and $n$ active faces,
coupled to a live block $B$ affine of order $m$ in the sense of
\cite[\S3.1]{paper}. Assume \hyplabel{H0$'$}--\hyplabel{H3$'$},
\hyplabel{H5$'$}, the live part of \hyplabel{H4$'$}, and \hyplabel{H3-gap}.
Then there is $\beta_0<\infty$ such that for $\beta\ge\beta_0$ all five clauses
\textup{(i)}, \textup{(ii)}, \textup{(iii)}, \textup{(ivA)}, \textup{(ivB)} of
\cite[Thm.~6]{paper} hold verbatim with
\[
c\ \longmapsto\ C_{i^\ast i^\ast},\qquad
R\ \longmapsto\ R_{i^\ast},\qquad
\delta_1\ \longmapsto\ \dgap ,
\]
release occurring on the face $i^\ast$. Moreover the threshold satisfies
\begin{equation}\label{eq:beta0A}
\beta_0\ =\ \max\Bigl(\beta_0^{\text{\cite{paper}}},\
\Theta\bigl(\dgap^{-1}\log(1/\dgap)\bigr)\Bigr),
\end{equation}
and, if the $n$ data lie in general position in a bounded set so that the
typical margin gap is $\dgap=\Theta(1/n)$, then $\beta_0=O(n\log n)$.
\end{theorem}

\begin{proof}
By \cref{lem:A1} the leader's quadrature is the identity of
\cite[L1$'$]{paper} with the stated substitutions and an error of the same
species as $\mathcal E$ (exponentially small in $\beta\dgap$ instead of
$\beta\delta_1$). Everything downstream of L1$'$ in \cite[App.~B]{paper} is then
literally unchanged: L2$'$ is a statement about $\theta_B$ and never sees the
faces of $A$; L3$'$ and L6$'$ are corner analyses in the single variable
$z_{i^\ast}$ --- and step~4 of L6$'$, the \emph{scalar reduction}, is the one
step that consumes the $\delta_1$ clause of \hyplabel{H4$'$}, which
\cref{lem:A1} has now supplied in the form ``all other faces contribute
$O(e^{-\beta\dgap})$''; L4$'$ is stated for prescribed $R$ and applies to
$R_{i^\ast}$; L5$'$ is the bootstrap and is arithmetic in the same objects.
\Cref{lem:A2} supplies the one genuinely new ingredient, which the scalar
theorem does not need and which is invisible in it: that no \emph{other} face
revives first. Release therefore occurs on $i^\ast$ at the first zero of
$\Ar_{i^\ast}$, up to the corrections of clause~(ii).

For \eqref{eq:beta0A}: \eqref{eq:H3gap} reads
$\beta\dgap\ge\log\beta+\log(\Gamma/\eta)$. Setting $\beta=\beta_0$ and solving
the self-consistent inequality, $\beta_0\dgap\ge\log\beta_0+K$ is satisfied as
soon as $\beta_0\ge2\dgap^{-1}\log(\dgap^{-1})$ for $\dgap$ small and $K$ fixed
(substitute: $\beta_0\dgap\ge2\log(1/\dgap)$ while
$\log\beta_0=\log(2/\dgap)+\log\log(1/\dgap)\le
\tfrac32\log(1/\dgap)$ for $\dgap$ below an absolute constant). Taking the max
against the paper's own five rows gives \eqref{eq:beta0A}; substituting
$\dgap=\Theta(1/n)$ gives $\beta_0=O(n\log n)$.
\end{proof}

\begin{remark}[what \cref{thm:A} buys, stated flatly]\label{rem:whatAbuys}
No new mathematics. Three things:
\emph{(1)} it identifies $\delta_1$ of \hyplabel{H4$'$} --- which
\cite[Rem.~hyp]{paper} treats as removable bookkeeping --- as the
\emph{essential and measurable} hypothesis of the multi-data problem, namely
leader slack;
\emph{(2)} it says that the governing parameter is $\beta\dgap$ and \textbf{not}
$\beta$, which is what the toy of \cite[\S4.1]{scoping} sees, the transition
sitting at $\beta\dgap\approx8$--$10$ stably over two decades of $\beta$;
\emph{(3)} it prices the hypothesis, $\beta_0=O(n\log n)$, which is a falsifiable
prediction about how the theorem degrades with the number of data. Against it
stands the measurement of \cite[\S3]{scoping}: in $9/9$ real episodes
$\beta\dgap\le0.35$. \textbf{The hypothesis of \cref{thm:A} is violated in every
real episode measured so far}; the scalar law nevertheless succeeds $13/16$
there, which is a fact about robustness and not about the hypothesis.
\end{remark}

\begin{remark}[the hypothesis is checkable, and it is not an assumption on $B$]
\hyplabel{H3-gap} involves $z^0$, $C$ and the forcing --- all read off a single
hybrid integration, exactly like \hyplabel{H1$'$}--\hyplabel{H2$'$}. It is
\emph{not} a condition on $m$ or on the live block. The one non-checkable
ingredient is the assumption inside \cref{lem:A1} that the leader remains the
leader on the window; \cref{lem:A2} makes that consistent (no other face can
reach zero) but does not prove it, since a face could in principle overtake the
leader while both remain parked. Closing that is
\cref{sec:open}\,(D1).
\end{remark}

%======================================================================
\section{Theorem B: the per-face diagonal law}\label{sec:B}

Here is the substantive statement. It abandons the leader entirely: no face is
privileged, and the release is decided by a \emph{race between $n$ clocks with
$n$ different thresholds}.

\subsection{$C$-orthogonality, and the tolerance}

\begin{itemize}
\item[\hyplabel{H7}] \textbf{$C$-orthogonality of the active faces.} With
\[
\varepsilon:=\max_{i\ne k}\frac{|C_{ik}|}{C_{ii}},\qquad
W:=\sup_{t\in[\thit,\trel]}\max_{i,k}\frac{1+u_i(t)}{1+u_k(t)}
\quad(\text{the weight of \cref{rem:weight}}),
\]
one asks $\varepsilon_{\mathrm{eff}}:=(n-1)\,\varepsilon\,W\le\varepsilon_0$
for a tolerance $\varepsilon_0$ fixed below.
\end{itemize}

The exactly-orthogonal case is $\varepsilon=0$ and then $\varepsilon_{\rm
eff}=0$ whatever $W$ does; it is realizable, since only $C$ enters
\eqref{eq:D1}: for $n=2$ take $x_1x_2=-1$, and in general $n\le d+1$ mutually
orthogonal augmented data. For $\varepsilon>0$ the product with $W$ is what
matters and $W$ can be exponentially large --- see \cref{rem:whyW}.

\begin{equation}\label{eq:tau}
\text{\textbf{Per-face threshold:}}\qquad
\tau_i:=\frac{\Psi(u_i^0)-1}{\beta C_{ii}},
\qquad
\Psi(u_i^0)=\beta\delta_i+e^{\beta\delta_i},
\end{equation}
so $\tau_{i^\ast}=0$ for the leader ($\delta=0$, $u^0=1$, $\Psi=1$) and
$\tau_i\approx e^{\beta\delta_i}/(\beta C_{ii})$ for a face parked at depth
$\delta_i$: \emph{the area a face needs is exponential in its depth.} Define
\begin{equation}\label{eq:ti}
t_i:=\inf\{t>\thit:\ \Ar_i(t)\ge\tau_i\},\qquad
i_{\rm win}:=\arg\min_it_i,\qquad t_{\rm win}:=\min_it_i .
\end{equation}

\subsection{The two lemmas}

\begin{lemma}[exact decoupling at $\varepsilon=0$]\label{lem:B1}
If $C_{ik}=0$ for all $i\ne k$ among the active faces, then \eqref{eq:D4}
becomes, for each $i$ separately and \emph{exactly},
\begin{equation}\label{eq:Bdec}
\frac{d\Psi(u_i)}{dt}=-\beta C_{ii}R_i(t)+C_{ii}\log(1+w_i),
\end{equation}
which is \cite[eq.~(quadratureF)]{paper} with $c=C_{ii}$, $R=R_i$ and
$\mathcal E\equiv0$. Consequently each face obeys \cite[L4$'$]{paper}
separately: $\Psi(u_i(t))=\Psi(u_i^0)-\beta C_{ii}\Ar_i(t)+{\rm corr}_i(t)$ with
$0\le{\rm corr}_i=\int C_{ii}\log(1+w_i)=O(\log\beta/\beta)$, and face $i$
returns to $z_i=0$ (i.e.\ $\Psi(u_i)=\Psi(1)=1$) at
$t_i+O(\log\beta/\beta^{2})$, with $t_i$ as in \eqref{eq:ti}.
\end{lemma}

\begin{proof}
Immediate from \eqref{eq:D4}: the coupling term is a sum over $k\ne i$ of terms
carrying the factor $C_{ik}=0$. There is no approximation and no $\delta_1$
bookkeeping to do, because there is nothing to bound --- and note that this is
\emph{stronger} than the scalar theorem's situation, where the other faces
contribute $O(e^{-\beta\delta_1})$ rather than $0$. The rest is
\cite[L4$'$(i)--(ii)]{paper} applied verbatim to the scalar equation
\eqref{eq:Bdec}, whose Chaplygin comparison is available because the right-hand
side is decreasing in $\Psi$; the threshold form \eqref{eq:tau} is the
statement $\Psi(u_i)=1$ solved for $\Ar_i$ with ${\rm corr}_i$ dropped, and the
$O(\log\beta/\beta^2)$ is ${\rm corr}_i$ divided by the transversal speed
$\beta C_{ii}R_i(t_i)$, i.e.\ \cite[L4$'$(ii)]{paper}. Verified numerically to
$1.3\cdot10^{-7}$ in \cite[I3]{scoping}.
\end{proof}

The new lemma is the next one. \Cref{lem:B1} gives $n$ clocks; nothing yet says
that the \emph{first} clock to ring is the one the exact flow obeys, nor that
the ordering of the $t_i$ survives the perturbation $\mathcal C_i$.

\begin{lemma}[stability of the $\arg\min$ across faces]\label{lem:argmin}
Suppose each face's exact $\Psi(u_i)$ differs from the decoupled solution of
\cref{lem:B1} by at most $\eta_i$ in sup-norm on $[\thit,t_{\rm win}+\Delta]$,
$\Delta>0$ fixed. Suppose:
\begin{itemize}
\item[\textup{(T)}] \emph{transversality of the winner}:
$R_{i_{\rm win}}(t_{\rm win})>0$ and $\Ar_{i_{\rm win}}$ crosses $\tau_{i_{\rm
win}}$ transversally, so that
$\frac{d}{dt}\Psi(u_{i_{\rm win}})\le-\tfrac12\beta C_{ii}R_{i_{\rm win}}(t_{\rm
win})=:-\beta\nu<0$ in a neighbourhood of $t_{\rm win}$;
\item[\textup{(S)}] \emph{threshold separation}: for every $k\ne i_{\rm win}$,
\[
\sigma_k:=\Psi\bigl(u_k(t_{\rm win})\bigr)-1
\ \ge\ \beta C_{kk}\bigl(\tau_k-\Ar_k(t_{\rm win})\bigr)-O(\log\beta)
\ >\ \eta_k .
\]
\end{itemize}
Then the exact system releases on the face $i_{\rm win}$, at a time
\begin{equation}\label{eq:argminshift}
\bigl|t^{\rm exact}_{\rm release}-t_{\rm win}\bigr|
\ \le\ \frac{\eta_{i_{\rm win}}}{\beta\nu}\ +\ O\!\left(\frac{\log\beta}{\beta^2}\right).
\end{equation}
\end{lemma}

\begin{proof}
Release of face $j$ means $\Psi(u_j)=1$ for the first time (since $\Psi$ is
decreasing on $u\ge1$ with $\Psi(1)=1$, and $z_j\ge0\iff u_j\le1\iff\Psi\le1$;
$\Psi$ is minimized at $u=1$ over $u\ge1$, so $\Psi(u_j)>1$ strictly while
$z_j<0$). By (S), at $t=t_{\rm win}$ every non-winner has
$\Psi(u_k)\ge1+\sigma_k$ in the decoupled solution, hence
$\ge1+\sigma_k-\eta_k>1$ in the exact one: \emph{no non-winner has released by
$t_{\rm win}$.} It remains to show that the winner does, within the stated
window. By \cref{lem:B1} the decoupled winner satisfies $\Psi=1$ at
$t_{\rm win}+O(\log\beta/\beta^2)$; by (T) it is strictly decreasing at rate
$\ge\beta\nu$ there, so the exact $\Psi$, which differs by at most
$\eta_{i_{\rm win}}$, attains $1$ at a time differing by at most
$\eta_{i_{\rm win}}/(\beta\nu)$ --- this is the elementary
``$|f^{-1}(y)-g^{-1}(y)|\le\|f-g\|_\infty/\inf|f'|$'' for the strictly monotone
$\Psi$ near its crossing. Adding the two displacements gives
\eqref{eq:argminshift}. Both (T) and (S) are needed and neither is implied by
the other: (T) alone permits a non-winner to have crossed earlier, (S) alone
permits the winner's crossing time to be arbitrarily sensitive.
\end{proof}

\begin{remark}[(S) is a hypothesis about the \emph{instance}, not about $\beta$]
$\sigma_k$ is the runner-up's remaining slack in the $\Psi$ coordinate at the
moment the winner rings; the failure mode it excludes is a \emph{tie}, and a tie
is exactly the degenerate configuration seen in the real episode seed~1923 of
\cite[\S3.2]{scoping}, where two faces sit at relative weight $0.997$ for the
whole episode. In that instance no per-face law of this kind can predict which
face returns, and this note does not claim otherwise. In the language of
\cite{scoping}, (S) says the trajectory is bounded away from the tropical wall.
\end{remark}

\begin{theorem}[per-face diagonal area law]\label{thm:B}
Let $A$ be a softplus-$\beta$ unit with readout $+1$ and $n$ active faces,
coupled to a live block $B$ affine of order $m$. Assume
\hyplabel{H0$'$}--\hyplabel{H3$'$}, \hyplabel{H5$'$}, the live part of
\hyplabel{H4$'$}, hypothesis \hyplabel{H7} with $\varepsilon=0$, and the
conditions \textup{(T)}, \textup{(S)} of \cref{lem:argmin}. Then there is
$\beta_0<\infty$ such that for $\beta\ge\beta_0$:
\begin{enumerate}
\item[(i)] \emph{(per-face sticking)} on compacts where $\Ar_i<\tau_i$, face $i$
parks at $u_i=\Psi^{-1}\bigl(\Psi(u_i^0)-\beta C_{ii}\Ar_i\bigr)(1+o(1))$; for
the leader this is $u=C_{i^\ast i^\ast}\beta|\Ar_{i^\ast}|(1+o(1))$, i.e.\
clause~(i) of \cite[Thm.~6]{paper};
\item[(ii)] \emph{(which face)} release occurs on $i_{\rm win}=\arg\min_it_i$
with $t_i$ the first time $\Ar_i$ reaches the threshold $\tau_i$ of
\eqref{eq:tau};
\item[(iii)] \emph{(when)} the release instant obeys
\begin{equation}\label{eq:relB}
t_{\rm release}
=t_{\rm win}
+\frac{{\rm corr}_{i_{\rm win}}(\beta)}
{\beta\,C_{i_{\rm win}i_{\rm win}}\,R_{i_{\rm win}}(t_{\rm win})}\,(1+o(1))
=t_{\rm win}+O\!\left(\frac{\log\beta}{\beta^2}\right);
\end{equation}
\item[(iv)] clauses \textup{(iii)}, \textup{(ivA)}, \textup{(ivB)} of
\cite[Thm.~6]{paper} hold with $c=C_{i_{\rm win}i_{\rm win}}$,
$R=R_{i_{\rm win}}$.
\end{enumerate}
If instead $\varepsilon>0$, the same conclusions hold with $t_{\rm win}$
displaced by \eqref{eq:argminshift} and
\begin{equation}\label{eq:etaeps}
\eta_i\ \le\ \beta\,\varepsilon_{\mathrm{eff}}\,C_{ii}\,(\Cres+1)\,\Tep ,
\end{equation}
provided $\eta_k<\sigma_k$ for every $k$ --- which is the quantitative meaning
of \hyplabel{H7}, with $\varepsilon_0$ read off as
$\varepsilon_0=\min_k\sigma_k/\bigl(\beta C_{kk}(\Cres+1)\Tep\bigr)$.
\end{theorem}

\begin{proof}
For $\varepsilon=0$: \cref{lem:B1} gives $n$ exact scalar problems of the
\cite[L1$'$/L4$'$]{paper} form, one per face, and \cref{lem:argmin} with
$\eta_i=0$ selects the winner and gives clause~(ii); clause~(iii) is
\cite[L4$'$(ii)]{paper} applied to the winner; clause~(i) is
\cite[L4$'$(i)]{paper} inverted, which for $u^0=1$ reduces to the paper's
$u=c\beta|\Ar|$ and for $u^0\gg1$ carries $\Psi(u^0)$; clause~(iv) is
L2$'$/L6$'$, which never see the faces of $A$ except through the scalar
reduction, and that reduction is now supplied by \cref{lem:argmin}(S) rather
than by $\delta_1$: at the release instant every other face is bounded away from
its own kink by $\sigma_k$.

For $\varepsilon>0$, bound $|\mathcal C_i|$ in \eqref{eq:D4} by
$\beta\sum_{k\ne i}|C_{ik}|\cdot W\cdot(\Cres+\max g_\beta)
\le\beta C_{ii}\varepsilon_{\mathrm{eff}}(\Cres+1)$ and integrate over the
episode, which is \eqref{eq:etaeps}; then apply \cref{lem:argmin}.
\end{proof}

\begin{remark}[why $W$ cannot be dropped from \hyplabel{H7}]\label{rem:whyW}
It is tempting to state $C$-orthogonality as ``$\varepsilon$ small'' alone. That
is false, and the falsity is the content of \cref{rem:weight}: a face parked at
depth $\delta_i$ against a face at depth $\delta_k\ll\delta_i$ sees the coupling
amplified by $e^{\beta(\delta_i-\delta_k)}$. What saves the statement is that
the \emph{winner} is, at its own crossing, the shallowest face, so $W\le1$ for
$i=i_{\rm win}$ at $t\approx t_{\rm win}$; but $\eta_k$ for the deep runners-up
is genuinely amplified, and it is bounded against $\sigma_k$ --- which is also
exponentially large, $\sigma_k\approx e^{\beta\delta_k}$. The two exponentials
are of the same species and the condition $\eta_k<\sigma_k$ is not vacuous. We
have \emph{not} extracted a clean sufficient condition in terms of
$(\varepsilon,\beta,\delta)$ alone; \cref{sec:open}\,(D3).
\end{remark}

\begin{remark}[interpretation: one clock per datum]
The physical reading is worth stating because it is what the multi-data law
\emph{adds} to the paper. Each datum carries its own area clock, and the datum
parked at depth $\delta$ needs an area larger by $e^{\beta\delta}$. That
exponential factor --- and not the hypothesis \hyplabel{H4$'$} --- is what makes
the leader win when $\beta$ is large; and it is exactly what \emph{fails} to be
large when $\beta\delta=O(1)$, which is the regime of every real episode
measured in \cite[\S3]{scoping}. \cref{thm:A} and \cref{thm:B} are therefore
not two theorems about different networks: they are the two sides of the single
quantity $\beta\delta$.
\end{remark}

%======================================================================
\section{Numerical verification, in the regime the scoping note could not
reach}\label{sec:num}

\subsection{Design, and the pre-registered killers}

The obstacle (A5) of \cite{scoping} is that the regime where the leader
hypothesis holds ($\beta\dgap\gtrsim10$) is the regime where the $n{=}16$ MLP
campaign observes no revivals at all. The design here attacks it directly, as
\cite[\S6]{scoping} recommends: $n\in\{2,3\}$ data and high $\beta$, so that
$\beta\dgap\ge10$ holds \emph{by construction} in the block that tests
\cref{thm:A}.

Since the field \eqref{eq:D1} sees the data only through $C$, we draw $C$
directly: diagonal (exact \hyplabel{H7}, $\varepsilon=0$) in blocks (a) and (b),
and $C=XX^\top/n$ with real $x_i\in[0.6,1.6]$ in block (c), which is strongly
coupled. The forcing is prescribed and affine, $R_i(t)=-a_i+b_it$ with
$a_i,b_i>0$ (so \hyplabel{H1$'$}--\hyplabel{H2$'$} hold per face) --- this is
the forcing class of \cite[Thm.~6]{paper}. Integration is RK4 on the exact
\eqref{eq:D1} with $dt=\min(10^{-4},1/(20\beta))$, zero-crossing interpolated
linearly; $\beta\in\{200,500,10^3,2\cdot10^3\}$; $48$ instances; a $dt/2$
control on every instance.

Predictors, both computed from the forcing alone:
\textbf{S} $=$ the scalar leader law of \cref{thm:A} ($\tau_{i^\ast}=0$, face
$i^\ast$, plus the first-order coefficient of \cite[Thm.~6(ii)]{paper});
\textbf{B} $=$ the per-face law \eqref{eq:ti} ($\min_i$, with the same
first-order coefficient at the winner). Conditions fixed before the run:
\emph{(a) kills \cref{thm:A}} if S predicts the wrong face or its error does not
decay in $\beta$; \emph{(b) kills \cref{thm:B}} likewise for B; \emph{(c)}
delimits: if B kept working at $\varepsilon\approx1.2$, \hyplabel{H7} would be
unnecessary and \cref{thm:B} misstated.

\subsection{Result}

\begin{center}
\begin{tabular}{@{}llrrrr@{}}
\toprule
block & regime & $N$ & face S & face B & $\varepsilon=\max_{i\ne k}|C_{ik}|/C_{ii}$\\
\midrule
(a) & $\beta\dgap\in[10.0,15.7]$, $\varepsilon=0$ & 18 & $18/18$ & $18/18$ & $0$\\
(b) & $\beta\dgap\in[0.57,4.0]$, $\varepsilon=0$  & 18 & $\mathbf{0/18}$ & $18/18$ & $0$\\
(c) & $\beta\dgap\in[0.71,3.6]$, coupled          & 12 & $12/12$ & $\mathbf{1/12}$ & $1.01$--$1.43$\\
\bottomrule
\end{tabular}
\end{center}

\begin{center}
\begin{tabular}{@{}llrrr@{}}
\toprule
block & predictor & median rel.\ error & median rel.\ error, zeroth order & $\log$--$\log$ slope in $\beta$\\
\midrule
(a) & S $\equiv$ B & $1.19\cdot10^{-3}$ & $1.51\cdot10^{-5}$ & $-0.71$\\
(b) & S & $2.44$ & $2.44$ & $-0.09$\\
(b) & B & $1.32\cdot10^{-3}$ & $2.00\cdot10^{-5}$ & $-1.23$\\
(c) & S & $0.30$ & $0.30$ & $-0.05$\\
(c) & B & $0.51$ & $0.51$ & $-0.07$\\
\bottomrule
\end{tabular}
\end{center}

\noindent
The $dt/2$ control gives a maximum relative displacement of $1.9\cdot10^{-7}$
over all $48$ instances, three to four orders below every error discussed, so
none of the numbers is an integration artefact. Reading:

\begin{itemize}
\item \textbf{(a) confirms \cref{thm:A}, including \cref{lem:A2}.} With
$\beta\dgap\ge10$ the two predictors \emph{coincide identically}, which is the
content of \cref{lem:A2}: the leader's threshold is $0$ and every other
threshold is $\tau_i\approx e^{\beta\delta_i}/(\beta C_{ii})$, in these
instances $3.7\cdot10^2$ to $2.1\cdot10^3$ in area units, unreachable. Face
correct $18/18$, error decaying with $\beta$.
\item \textbf{(b) is the discriminating block, and it discriminates.} With
$\beta\dgap\in[0.57,4]$ the scalar leader law predicts the \textbf{wrong face in
$18/18$ instances} and a time error of $244\%$ that does not decay
($\beta^{-0.09}$); the per-face law is right on $18/18$ faces with median error
$1.3\cdot10^{-3}$ decaying as $\beta^{-1.2}$. This reproduces, in a regime three
orders of magnitude higher in $\beta$, the qualitative failure that
\cite[\S4.1]{scoping} measured at $\beta\le10^3$ and that the real campaign sees
as $3/16$ episodes without area.
\item \textbf{(c) delimits \cref{thm:B}, as pre-registered.} At
$\varepsilon\approx1.2$ --- off-diagonal \emph{larger} than diagonal, the regime
of \cite[\S4.2]{scoping} where pushing face $2$ raises face $1$ faster than
itself --- the diagonal law predicts the wrong face in $11/12$ instances with a
$51\%$ error that does not decay ($\beta^{-0.07}$). \hyplabel{H7} is therefore
not decorative. Note that in this block the \emph{scalar} law gets the face
right $12/12$ while still missing the time by $30\%$: with strong positive
coupling everything is pushed towards the shallowest face, so the face is easy
and the time is not. Neither law is valid here, and the correct object is the
reduced flow (D6) of \cite{scoping}.
\end{itemize}

\subsection{The one number that does not close}\label{sec:naoFecha}

The tables above report the zeroth-order error (predictor $=t_{\rm win}$)
separately, and it is \emph{smaller than the first-order one by a factor
$\approx80$}: $1.5\cdot10^{-5}$ against $1.19\cdot10^{-3}$ in block (a),
$2.0\cdot10^{-5}$ against $1.32\cdot10^{-3}$ in block (b), with the first-order
predictor worse in $36/36$ instances. The measured signed displacement is
\[
\beta\bigl(t_{\rm release}-t_{\rm win}\bigr)=+0.004\ \text{to}\ +0.12,
\qquad
\frac{\beta^2\bigl(t_{\rm release}-t_{\rm win}\bigr)}{\log\beta}=0.6\ \text{to}\ 9.2,
\]
i.e.\ $t_{\rm release}-t_{\rm win}=+\Theta(\log\beta/\beta^2)$, \emph{positive},
of the order and sign of \eqref{eq:relB} --- which is \cite[L4$'$(ii)]{paper}.
It is \textbf{not} $-1/(cR^{\hyb}(\texit))\cdot\beta^{-1}$, the coefficient of
\cite[Thm.~6(ii)]{paper}, which in these instances is $-0.3$ to $-3.1$ in units
of $\beta^{-1}$ and would be $50$--$100\times$ larger than what is observed, with
the opposite sign.

A control was run to check whether this is an artefact of the multi-face
setting: $n=1$, $c=1$, $R=-1+1.2t$, the literal scalar scenario of the paper, at
$\beta\in\{2\cdot10^2,10^3,5\cdot10^3,2\cdot10^4\}$. Result:
$\beta(\trel-\texit)=+0.0625,\,+0.0157,\,+0.0037,\,+0.0010$ against the paper's
predicted $-1/(cR)=-1.000$, while
$\beta^2(\trel-\texit)/\log\beta=2.36,\,2.27,\,2.19,\,2.06$ --- flat, i.e.\
$\trel-\texit=\Theta(\log\beta/\beta^2)$ in the scalar case too.

\textbf{We state this without softening: in this toy, in the scalar case as much
as in the multi-face case, the first-order coefficient of
\cite[Thm.~6(ii)]{paper} is not observed, and what is observed is the
$O(\log\beta/\beta^2)$ displacement of \cite[L4$'$(ii)]{paper}.} Those two
statements of \cite{paper} are mutually inconsistent at order $\beta^{-1}$, and
this note \emph{cannot} adjudicate between them: the paper's own campaigns
measure the coefficient inside the \emph{coupled} flow, where $\texit$ is the
zero of the hybrid area of a live block that itself moves, and our forcing is
prescribed, so the two setups may not be measuring the same displacement. What
is safe, and is what \cref{thm:B}(iii) claims, is the prescribed-forcing
statement \eqref{eq:relB}. We do not edit the paper on the strength of a toy;
see \cref{sec:open}\,(D5).

\begin{remark}[executable]
Script \path{05-experimentos/t2-n-neuronios/degrau3a_verifica.py}, output
\path{degrau3a_verifica.jsonl} (one line per instance, plus four lines
\texttt{bloco="n1"} for the scalar control). Two workers, \texttt{nice -19}.
Every number in \cref{sec:num} is a field of that file.
\end{remark}

%======================================================================
\section{Open steps, declared}\label{sec:open}

\begin{itemize}
\item[(D1)] \textbf{The leader is assumed to stay the leader.} \Cref{lem:A1}
assumes $i^\ast$ constant on the window; \cref{lem:A2} shows no other face can
\emph{revive} first, which is weaker. A face could overtake the leader while
both are parked, and then \eqref{eq:Abound} loses its hypothesis mid-window.
Closing this needs a continuous-induction argument of the
\cite[L5$'$]{paper} type on the gap $z_{i^\ast}-z_{(2)}$, and it is not written.
\item[(D2)] \textbf{The leader-swap corner is outside this draft.} The genuinely
new mathematics of Degrau~3, per \cite[\S6]{scoping}, is the layer analysis at
the tropical wall $z_i+\beta^{-1}\log|R_i|=z_k+\beta^{-1}\log|R_k|$, where two
channels have comparable weight and the scalar reduction of
\cite[L6$'$]{paper} fails --- a two-fold singular perturbation in dimension
$d+1$. Hypothesis (S) of \cref{lem:argmin} is exactly the assumption that the
trajectory stays off that wall. Nothing here touches the corner.
\item[(D3)] \textbf{\hyplabel{H7} is not reduced to a clean inequality.} The
tolerance is stated as $\eta_k<\sigma_k$ (\cref{thm:B}, \cref{rem:whyW}), which
involves two exponentially large quantities of the same species. A sufficient
condition in terms of $(\varepsilon,\beta,\delta_i)$ alone is not extracted.
The numerics bracket the answer only crudely: $\varepsilon=0$ works,
$\varepsilon\approx1.2$ fails, and nothing between the two was drawn.
\item[(D4)] \textbf{The coupled case has no theorem.} For $\varepsilon=\Theta(1)$
there is no closed quadrature at all, and the reason is a count, not a
technicality: with $n>d+1$ active faces the condition for a $\Psi$ with
$d\Psi/dt$ a function of $t$ alone is $\nabla\Psi\cdot X_i=\kappa_ie^{-\beta
z_i}$, $n$ equations for $d+1$ unknowns, generically insoluble
\cite[\S2.2]{scoping}. The substitute is the reduced flow (D6) --- $d+1$
dimensions, still driven by the forcing alone, hence still in the
\cite[Thm.~6]{paper} family --- and block~(c) above is consistent with it, but
this note proves nothing about it.
\item[(D5)] \textbf{The first-order coefficient does not reproduce.} See
\cref{sec:naoFecha}. This is an inconsistency \emph{inside} \cite{paper}
(Thm.~6(ii) versus L4$'$(ii)) surfaced by a toy that sides with L4$'$(ii),
including in an $n=1$ control. It is not adjudicated here and the paper is not
edited. The cheapest next test is to rerun the paper's own
\path{verifica_forcing.py} instances and report
$\beta(\trel-\texit)$ against both candidates.
\item[(D6)] \textbf{The bridge to real data is not built (obstacle A5).} The
campaign here is a toy with prescribed affine forcing and $C$ drawn by hand. In
all $9$ real episodes instrumented in \cite[\S3]{scoping},
$\beta\dgap\le0.35$, so \emph{neither} \cref{thm:A} (which needs
$\beta\dgap\gtrsim10$) \emph{nor} the real network's $C=XX^\top/n$ (which is far
from diagonal) is in scope. What this note establishes is that both laws are
correct \emph{in their stated regimes} and that the regimes are separated by a
sharp, measurable quantity. Whether a real training run can be placed in either
regime is untested, and would require the pre-registered $n=2$--$3$, high-$\beta$
design of \cite[\S6(ii)]{scoping}.
\item[(D7)] \textbf{$N$ is small and the transition is not located.} $18/18/12$
instances, four values of $\beta$, one forcing family (affine, one sign change).
Block~(b) deliberately biases the leader's forcing to be slow, so it tests the
two sides of the transition but does \emph{not} measure where it sits; the
$\beta\dgap\approx8$--$10$ of \cite[\S4.1]{scoping} is the only estimate we
have, and \eqref{eq:H3gap} predicts $\log(\beta\Gamma)$, which at these $\beta$
is $5$--$8$ --- the same order, which is agreement of magnitude and not a fit.
\end{itemize}

\bigskip
\noindent\textbf{Etiquetas epist\^emicas.}
\Cref{eq:D1} and \eqref{eq:D4}: \emph{[ESTABELECIDO]} (identities, derived and
verified numerically in \cite{scoping}, errors $6\cdot10^{-11}$ and
$9.6\cdot10^{-9}$).
\Cref{lem:B1} at $\varepsilon=0$: \emph{[ESTABELECIDO]} (exact decoupling; it is
$n$ copies of \cite[L1$'$]{paper}).
\Cref{lem:A1,lem:A2,lem:argmin}: \emph{[ESPECULA\c{C}\~AO com deriva\c{c}\~ao]}
--- the computations are done and elementary, no step audited by a referee.
\Cref{thm:A,thm:B} as wholes: \emph{[ESPECULA\c{C}\~AO]} --- proof at lemma
granularity, resting on \cite{paper} App.~B for everything downstream of L1$'$,
with (D1) open.
Numbers of \cref{sec:num}: \emph{[RESULTADO 2026-07-27]}.
The failure of the first-order coefficient (\cref{sec:naoFecha}):
\emph{[RESULTADO 2026-07-27]} for the measurement,
\emph{[N\~AO VERIFICADO]} for its cause.

\begin{thebibliography}{9}
\bibitem{paper} L.~Caloi Santos, \emph{An entry--exit law for a dying softplus
unit} (\texttt{06-notas/paper/paper-v4.tex}), 2026. Theorem~6 (forcing form),
\S3.1; App.~B (proof of the six lemmas L1$'$--L6$'$); App.~B.9 (uniformity);
App.~B.10 (open items).
\bibitem{scoping} \emph{Degrau 3 --- scoping matem\'atico do teorema
multi-dado}, \texttt{06-notas/2026-07-27-degrau3-scoping.md}, 2026-07-27.
Identities (D1)--(D7), toy of \S4, replay of the real episodes in \S3.
\bibitem{degrau2} \emph{The Degrau~2 theorem: the entry--exit law with a
trainable readout}, \texttt{06-notas/t2-prova/teorema-degrau2-draft.tex},
2026-07-27.
\end{thebibliography}

\end{document}
